% ============================================================
% AUTOMATISCH GEGENEREERD BESTAND
% Niet handmatig aanpassen.
% ============================================================

\ifdefined\HCode
  % Geen afzonderlijke module-openingspagina in HTML.
\else
  \FVDmoduleopening
    {Wiskundige structuren en redeneringen}
    {In deze module verdiepen we ons in de logische en structurele fundamenten van de wiskunde. We starten met de taal van de logica. We onderzoeken wiskundige uitspraken en hun waarheidswaarde en leren werken met logische operatoren, implicaties en equivalenties. Met logische wetten en kwantoren leren we samengestelde uitspraken analyseren, correct formuleren en negeren. Daarbij besteden we bijzondere aandacht aan de opbouw en geldigheid van wiskundige redeneringen. Vervolgens richten we onze aandacht op de gehele getallen. We onderzoeken deelbaarheid, de Euclidische deling, priemgetallen en priemfactoren en ontdekken hoe modulorekenen toelaat om op een efficiënte manier over gehele getallen en hun eigenschappen te redeneren. Vanuit deze concrete eigenschappen van getallen en bewerkingen maken we ten slotte de stap naar algebraïsche structuren. We onderzoeken welke eigenschappen bewerkingen gemeenschappelijk hebben en hoe deze aanleiding geven tot structuren zoals groepen, ringen en lichamen. Zo maken we kennis met de kracht van abstractie: verschillende wiskundige objecten kunnen volgens dezelfde onderliggende structuur worden bestudeerd. Doorheen de module besteden we bijzondere aandacht aan wiskundig redeneren, argumenteren en bewijzen. We leren uitspraken kritisch onderzoeken, tegenvoorbeelden zoeken en redeneringen zorgvuldig opbouwen en beoordelen. Deze vaardigheden worden in de verdere cursus telkens opnieuw gebruikt en verder ontwikkeld. Logica, getaltheorie en algebraïsche structuren laten zien hoe wiskunde verder gaat dan rekenen alleen. Ze maken zichtbaar hoe wiskundigen vanuit precieze taal, geldige redeneringen en herkenbare structuren nieuwe eigenschappen ontdekken en algemene resultaten opbouwen.}
    {%
    \FVDmodulethemetile{1}{Logica — De taal van correct redeneren}
    \FVDmodulethemetile{2}{Deelbaarheid en getaltheorie — Structuur in de gehele getallen}
    \FVDmodulethemetile{3}{Algebraïsche structuren — Van concrete bewerkingen naar abstracte structuren}
    }
\fi
